\futureoutlook
\label{sec:future}

\subsection{Outstanding Tasks}

The pipeline architecture, the DFT reference, the MEAM merge logic, and a
composition-only sanity-check surrogate are all in place. The remaining
work concentrates on closing the loop with a parameter-level
neural-network optimization stage and on formal validation:

\begin{enumerate}
	\item \textbf{Implement the Stage~5 parameter optimization.} The
	      neural-network surrogate that maps MEAM parameters to alloy
	      elastic properties --- and the inverse-optimization step
	      that uses it to refit the Stage~4 library against
	      LAMMPS-evaluated targets --- is the headline remaining
	      task. The supporting infrastructure (parallel LAMMPS
	      sampling, MPI distribution, MEAM file I/O) is already in
	      place; what remains is to design a stable surrogate
	      architecture, choose a robust loss, and add an explicit
	      physical-stability penalty
	      ($C_{11}+C_{12}>0$, $C_{11}\geq C_{12}$, $\nu<0.48$) that
	      keeps the inverse optimizer inside the surrogate's domain
	      of validity.

	\item \textbf{Improve the DFT reference for Mn.} The
	      $\alpha$-Mn EOS scan does not converge and the Mn entry is
	      a tabulated fallback. As an interim measure a helper
	      backfills any missing or unphysical $C_{ij}$ from the
	      recorded $B$ and a per-element Poisson-ratio default, so
	      Stage~4 receives a complete and self-consistent set; this
	      also rescues Si, Ti, Zn and Mg whose non-cubic lattices skip
	      the elastic-constant computation entirely. A proper Mn
	      calculation (simpler reference structure or different
	      smearing scheme) is still planned.

	\item \textbf{Build a validation set of named alloys.} A formal
	      benchmark against industrial aluminium alloys (Al~2024,
	      Al~2219, Al~5052, Al~7075, Al~6061, Al~6063, Al~5083) and
	      at least two non-aluminium alloys (a steel and a
	      Ti~alloy) will be assembled once the Stage~5 optimizer
	      starts producing physical results. Both LAMMPS-with-merged-MEAM
	      baselines and final-optimized predictions will be reported.

	\item \textbf{Write the final report and prepare the poster.}
	      Sections from the present interim report will be carried over
	      to the final report, with the Stage~5 results and the formal
	      validation added. The report tooling has been automated:
	      every figure and LaTeX table in the present document is
	      regenerated automatically from the pipeline outputs, so
	      refreshed LAMMPS or DFT data flow through to the report on
	      the next compile.
\end{enumerate}

The infrastructure required to execute these tasks at scale is already in
place. The pipeline has built-in checkpoint/resume support: each of the
DFT and MEAM-initialization stages writes its output to a deterministic
location and the orchestrator detects which of those outputs already
exist, so a re-run after a failure or a parameter change skips work
that has already been completed.

\subsection{Updated Semester Plan}

The Gantt chart of Figure~\ref{fig:gantt} reflects the work that has
already been completed and the schedule of the remaining tasks. The first
five work packages from the proposal (literature review, DFT reference
calculations, MEAM initialization, NN surrogate training, and the first
pass at inverse optimization) are now complete; validation, scaling, and
report/poster preparation occupy the rest of the semester.

\begin{figure}[h]
	\centering
	\begin{semestergantt}
		\workpackage{Literature Review}{1}{3}
		\workpackage{DFT Reference Calculations}{2}{5}
		\workpackage{MEAM Initialization Pipeline}{4}{7}
		\workpackage{NN Surrogate Training}{6}{9}
		\workpackage{Inverse Optimization}{8}{13}
		\workpackage{Validation \& Testing}{10}{13}
		\workpackage{Writing \& Documentation}{11}{15}
		\workpackage{Poster Preparation}{13}{15}
	\end{semestergantt}
	\caption{Semester plan with the actual progress to date. Items 1--4
	are complete; items 5--8 are the focus of the second half of the
	semester.}
	\label{fig:gantt}
\end{figure}

\subsection{Description of Work Packages}
\begin{itemize}
	\item \textbf{Literature Review (W1--W3).} Survey of MEAM, EAM and
	      neural-network potential literature. Completed.
	\item \textbf{DFT Reference Calculations (W2--W5).} Build of
	      Quantum~Espresso, pseudopotential management and the
	      EOS/elastic-constant/binary-pair pipeline. Completed.
	\item \textbf{MEAM Initialization Pipeline (W4--W7).} Auto-discovery
	      and merging of available MEAM library files, DFT-to-MEAM
	      analytic mapping, generation of the unified twelve-element
	      library. Completed.
	\item \textbf{NN Surrogate Training (W6--W9).} Sampling of the MEAM
	      parameter space, parallel LAMMPS evaluation and training of
	      the parameter--property surrogate. First pass complete; will
	      be re-run with a larger sample set.
	\item \textbf{Inverse Optimization (W8--W13).} Gradient-based
	      back-propagation through the surrogate to fit elastic
	      targets. First pass complete with the limitations described
	      above; ongoing.
	\item \textbf{Validation \& Testing (W10--W13).} Comparison against
	      named industrial alloys and external validation sets. Begun.
	\item \textbf{Writing \& Documentation (W11--W15).} Final report
	      and code documentation.
	\item \textbf{Poster Preparation (W13--W15).} Final poster for the
	      end-of-semester presentation.
\end{itemize}
